\subsection{\label{sec.gf.prod}Infinite products}

Let us now extend our FPS playground somewhat. We have made sense of infinite
sums. What about infinite products?

\subsubsection{\label{subsec.gf.prod.exa}An example}

We start with a \textbf{motivating example} (due to Euler, in
\cite[\S 328--329]{Euler48}), which we shall first discuss informally.

Assume (for the time being) that the infinite product%
\begin{equation}
\prod_{i\in\mathbb{N}}\left(  1+x^{2^{i}}\right)  =\left(  1+x^{1}\right)
\left(  1+x^{2}\right)  \left(  1+x^{4}\right)  \left(  1+x^{8}\right)
\cdots\label{eq.fps.prod.binary.1}%
\end{equation}
in the ring $K\left[  \left[  x\right]  \right]  $ is meaningful, and that
such products behave as nicely as finite products. Can we simplify this product?

We can observe that each $i\in\mathbb{N}$ satisfies $1+x^{2^{i}}%
=\dfrac{1-x^{2^{i+1}}}{1-x^{2^{i}}}$ (since $1-x^{2^{i+1}}=1-\left(  x^{2^{i}%
}\right)  ^{2}=\left(  1-x^{2^{i}}\right)  \left(  1+x^{2^{i}}\right)  $).
Multiplying these equalities over all $i\in\mathbb{N}$, we obtain%
\[
\prod_{i\in\mathbb{N}}\left(  1+x^{2^{i}}\right)  =\prod_{i\in\mathbb{N}%
}\dfrac{1-x^{2^{i+1}}}{1-x^{2^{i}}}=\dfrac{1-x^{2}}{1-x^{1}}\cdot
\dfrac{1-x^{4}}{1-x^{2}}\cdot\dfrac{1-x^{8}}{1-x^{4}}\cdot\dfrac{1-x^{16}%
}{1-x^{8}}\cdot\cdots.
\]
The product on the right hand side here is a \emph{telescoping product} --
meaning that each numerator is cancelled by the denominator of the following
fraction. Assuming (somewhat plausibly, but far from rigorously) that we are
allowed to cancel infinitely many factors from an infinite product, we thus
end up with a single $1-x^{1}$ factor in the denominator. That is, our product
simplifies to $\dfrac{1}{1-x^{1}}$. Thus, we obtain%
\begin{equation}
\prod_{i\in\mathbb{N}}\left(  1+x^{2^{i}}\right)  =\dfrac{1}{1-x^{1}}%
=\dfrac{1}{1-x}=1+x+x^{2}+x^{3}+\cdots. \label{eq.fps.prod.binary.2}%
\end{equation}


This was not very rigorous, so let us try to compute the product $\prod
_{i\in\mathbb{N}}\left(  1+x^{2^{i}}\right)  $ in a different way. Namely, we
recall a simple fact about finite products: If $a_{0},a_{1},\ldots,a_{m}$ are
finitely many elements of a commutative ring, then the product
\begin{equation}
\prod_{i=0}^{m}\left(  1+a_{i}\right)  =\left(  1+a_{0}\right)  \left(
1+a_{1}\right)  \cdots\left(  1+a_{m}\right)
\label{eq.fps.prod.binary.prod-fin}%
\end{equation}
equals the sum\footnote{The indices $i_{1},i_{2},\ldots,i_{k}$ of the sum are
supposed to be nonnegative integers.}
\[
\sum_{i_{1}<i_{2}<\cdots<i_{k}\leq m}a_{i_{1}}a_{i_{2}}\cdots a_{i_{k}}%
\]
of all the $2^{m+1}$ \textquotedblleft sub-products\textquotedblright\ of the
product $a_{0}a_{1}\cdots a_{m}$ (because this sum is what we obtain if we
expand $\prod_{i=0}^{m}\left(  1+a_{i}\right)  $ by repeatedly applying
distributivity). For example, for $m=2$, this is saying that%
\begin{align*}
&  \left(  1+a_{0}\right)  \left(  1+a_{1}\right)  \left(  1+a_{2}\right) \\
&  =1+a_{0}+a_{1}+a_{2}+a_{0}a_{1}+a_{0}a_{2}+a_{1}a_{2}+a_{0}a_{1}a_{2}.
\end{align*}
Now, it is plausible to expect the same formula $\prod_{i=0}^{m}\left(
1+a_{i}\right)  =\sum_{I\subseteq\left\{  0,1,\ldots,m\right\}  }%
\ \ \prod_{i\in I}a_{i}$ to hold if \textquotedblleft$m$ is $\infty
$\textquotedblright\ (that is, if the product ranges over all $i\in\mathbb{N}%
$), provided that the product is meaningful. In other words, it is plausible
to expect that%
\begin{equation}
\prod_{i\in\mathbb{N}}\left(  1+a_{i}\right)  =\sum_{i_{1}<i_{2}<\cdots<i_{k}%
}a_{i_{1}}a_{i_{2}}\cdots a_{i_{k}} \label{eq.fps.prod.binary.prod-inf}%
\end{equation}
for any infinite sequence $a_{0},a_{1},a_{2},\ldots$ as long as $\prod
_{i\in\mathbb{N}}\left(  1+a_{i}\right)  $ makes sense. (It's a little bit
more complicated than that, but we aren't trying to be fully rigorous yet. The
correct condition is that the sequence $\left(  a_{0},a_{1},a_{2}%
,\ldots\right)  $ is summable.) If we now apply
(\ref{eq.fps.prod.binary.prod-inf}) to $a_{i}=x^{2^{i}}$, then we obtain%
\begin{align}
\prod_{i\in\mathbb{N}}\left(  1+x^{2^{i}}\right)   &  =\sum_{i_{1}%
<i_{2}<\cdots<i_{k}}x^{2^{i_{1}}}x^{2^{i_{2}}}\cdots x^{2^{i_{k}}}=\sum
_{i_{1}<i_{2}<\cdots<i_{k}}x^{2^{i_{1}}+2^{i_{2}}+\cdots+2^{i_{k}}}\nonumber\\
&  =\sum_{n\in\mathbb{N}}q_{n}x^{n}, \label{eq.fps.prod.binary.5}%
\end{align}
where $q_{n}$ is the \# of ways to write the integer $n$ as a sum $2^{i_{1}%
}+2^{i_{2}}+\cdots+2^{i_{k}}$ with nonnegative integers $i_{1},i_{2}%
,\ldots,i_{k}$ satisfying $i_{1}<i_{2}<\cdots<i_{k}$. Comparing this with
(\ref{eq.fps.prod.binary.2}), we obtain%
\[
\sum_{n\in\mathbb{N}}q_{n}x^{n}=1+x+x^{2}+x^{3}+\cdots=\sum_{n\in\mathbb{N}%
}x^{n},
\]
at least if our assumptions were valid. Comparing coefficients, this would
mean that $q_{n}=1$ for each $n\in\mathbb{N}$. In other words, each
$n\in\mathbb{N}$ can be written in \textbf{exactly one} way as a sum
$2^{i_{1}}+2^{i_{2}}+\cdots+2^{i_{k}}$ with nonnegative integers $i_{1}%
,i_{2},\ldots,i_{k}$ satisfying $i_{1}<i_{2}<\cdots<i_{k}$. In other words,
each $n\in\mathbb{N}$ can be written uniquely as a finite sum of distinct
powers of $2$.

Is this true? Yes, because this is just saying that each $n\in\mathbb{N}$ has
a unique binary representation. For example, $21=2^{4}+2^{2}+2^{0}$
corresponds to the binary representation $21=\left(  10101\right)  _{2}$.

Thus, the two results we have obtained in (\ref{eq.fps.prod.binary.2}) and
(\ref{eq.fps.prod.binary.5}) are actually equal, which is reassuring. Yet,
this does not replace a formal definition of infinite products that rigorously
justifies the above arguments.

\subsubsection{\label{subsec.gf.prod.def}A rigorous definition}

One way of rigorously defining infinite products of FPSs can be found in
\cite[\S 7.5]{Loehr-BC}. However, this definition only defines infinite
products of the form $\prod_{i\in\mathbb{N}}$ or $\prod_{i=1}^{\infty}$, but
not (for example) of the form $\prod_{I\subseteq\mathbb{N}}$ or $\prod
_{\left(  i,j\right)  \in\mathbb{N}\times\mathbb{N}}$. Another definition of
infinite products uses the $\operatorname*{Log}$ and $\operatorname*{Exp}$
bijections from Definition \ref{def.fps.Exp-Log-maps} to turn products into
sums; but this requires $K$ to be a $\mathbb{Q}$-algebra (since
$\operatorname*{Log}$ and $\operatorname*{Exp}$ aren't defined otherwise).
Thus, we shall give a different definition here.

We recall our definition of infinite sums of FPSs (Definition
\ref{def.fps.summable}):

\begin{statement}
\textbf{Definition \ref{def.fps.summable} (repeated).} A (possibly infinite)
family $\left(  \mathbf{a}_{i}\right)  _{i\in I}$ of FPSs is said to be
\emph{summable} if%
\[
\text{for each }n\in\mathbb{N}\text{, all but finitely many }i\in I\text{
satisfy }\left[  x^{n}\right]  \mathbf{a}_{i}=0.
\]
In this case, the sum $\sum_{i\in I}\mathbf{a}_{i}$ is defined to be the FPS
with%
\[
\left[  x^{n}\right]  \left(  \sum_{i\in I}\mathbf{a}_{i}\right)
=\underbrace{\sum_{i\in I}\left[  x^{n}\right]  \mathbf{a}_{i}}%
_{\substack{\text{an essentially}\\\text{finite sum}}%
}\ \ \ \ \ \ \ \ \ \ \text{for all }n\in\mathbb{N}\text{.}%
\]

\end{statement}

This is how we defined infinite sums of FPSs. We cannot use the same
definition for infinite products, because usually%
\[
\text{we \textbf{don't} expect to have }\left[  x^{n}\right]  \left(
\prod_{i\in I}\mathbf{a}_{i}\right)  =\prod_{i\in I}\left[  x^{n}\right]
\mathbf{a}_{i}%
\]
(after all, multiplication of FPSs is not defined coefficientwise). The
condition \textquotedblleft all but finitely many $i\in I$ satisfy $\left[
x^{n}\right]  \mathbf{a}_{i}=0$\textquotedblright\ is therefore not what we
are looking for.

Let us instead go back to the idea behind Definition \ref{def.fps.summable}.
Let us fix some $n\in\mathbb{N}$. What was the actual purpose of the
\textquotedblleft all but finitely many $i\in I$ satisfy $\left[
x^{n}\right]  \mathbf{a}_{i}=0$\textquotedblright\ condition? The purpose was
to ensure that the coefficient $\left[  x^{n}\right]  \left(  \sum_{i\in
I}\mathbf{a}_{i}\right)  $ is determined by \textbf{finitely many} of the
$\mathbf{a}_{i}$'s. In other words, the purpose was to ensure that there is a
\textbf{finite} partial sum of $\sum_{i\in I}\mathbf{a}_{i}$ such that if we
add any further $\mathbf{a}_{i}$'s to this partial sum, then the coefficient
of $x^{n}$ does not change any more. Here is a way to restate this condition
more rigorously: There is a \textbf{finite} subset $M$ of $I$ such that every
finite subset $J$ of $I$ satisfying $M\subseteq J\subseteq I$ satisfies%
\[
\left[  x^{n}\right]  \left(  \sum_{i\in M}\mathbf{a}_{i}\right)  =\left[
x^{n}\right]  \left(  \sum_{i\in J}\mathbf{a}_{i}\right)  .
\]
(The subset $M$ here is the indexing set of our finite partial sum, and the
set $J$ is what it becomes if we add some further $\mathbf{a}_{i}$'s to this
partial sum.)

This condition is a mouthful; this is why we found it easier to boil it down
to the simple \textquotedblleft all but finitely many $i\in I$ satisfy
$\left[  x^{n}\right]  \mathbf{a}_{i}=0$\textquotedblright\ condition in the
case of infinite sums. However, in the case of infinite products, we cannot
boil it down to something this simple; thus, we have to live with it.

Fortunately, we can simplify our life by giving this condition a name. To
highlight the analogy, let us define it both for sums and for products:

\begin{definition}
\label{def.fps.determines-xn-coeff}Let $\left(  \mathbf{a}_{i}\right)  _{i\in
I}\in K\left[  \left[  x\right]  \right]  ^{I}$ be a (possibly infinite)
family of FPSs. Let $n\in\mathbb{N}$. Let $M$ be a finite subset of $I$.
\medskip

\textbf{(a)} We say that $M$ \emph{determines the }$x^{n}$\emph{-coefficient
in the sum of }$\left(  \mathbf{a}_{i}\right)  _{i\in I}$ if every finite
subset $J$ of $I$ satisfying $M\subseteq J\subseteq I$ satisfies%
\[
\left[  x^{n}\right]  \left(  \sum_{i\in J}\mathbf{a}_{i}\right)  =\left[
x^{n}\right]  \left(  \sum_{i\in M}\mathbf{a}_{i}\right)  .
\]
(You can think of this condition as saying \textquotedblleft If you add any
further $\mathbf{a}_{i}$s to the sum $\sum_{i\in M}\mathbf{a}_{i}$, then the
$x^{n}$-coefficient stays unchanged\textquotedblright, or, more informally:
\textquotedblleft If you want to know the $x^{n}$-coefficient of $\sum_{i\in
I}\mathbf{a}_{i}$, it suffices to take the partial sum over all $i\in
M$\textquotedblright.) \medskip

\textbf{(b)} We say that $M$ \emph{determines the }$x^{n}$\emph{-coefficient
in the product of }$\left(  \mathbf{a}_{i}\right)  _{i\in I}$ if every finite
subset $J$ of $I$ satisfying $M\subseteq J\subseteq I$ satisfies%
\[
\left[  x^{n}\right]  \left(  \prod_{i\in J}\mathbf{a}_{i}\right)  =\left[
x^{n}\right]  \left(  \prod_{i\in M}\mathbf{a}_{i}\right)  .
\]
(You can think of this condition as saying \textquotedblleft If you multiply
any further $\mathbf{a}_{i}$s to the product $\prod_{i\in M}\mathbf{a}_{i}$,
then the $x^{n}$-coefficient stays unchanged\textquotedblright, or, more
informally: \textquotedblleft If you want to know the $x^{n}$-coefficient of
$\prod_{i\in I}\mathbf{a}_{i}$, it suffices to take the partial product over
all $i\in M$\textquotedblright.)
\end{definition}

\begin{example}
\textbf{(a)} Consider the family
\begin{align*}
&  \left(  \left(  x+x^{2}\right)  ^{i}\right)  _{i\in\mathbb{N}}\\
&  =\left(  \left(  x+x^{2}\right)  ^{0},\ \ \left(  x+x^{2}\right)
^{1},\ \ \left(  x+x^{2}\right)  ^{2},\ \ \left(  x+x^{2}\right)
^{3},\ \ \left(  x+x^{2}\right)  ^{4},\ \ \ldots\right) \\
&  =\left(  1,\ \ x+x^{2},\ \ x^{2}+2x^{3}+x^{4},\ \ x^{3}+3x^{4}+3x^{5}%
+x^{6},\ \ x^{4}+4x^{5}+6x^{6}+4x^{7}+x^{8},\ \ \ldots\right)
\end{align*}
of FPSs. The subset $\left\{  2,3\right\}  $ of $\mathbb{N}$ determines the
$x^{3}$-coefficient in the sum of this family $\left(  \left(  x+x^{2}\right)
^{i}\right)  _{i\in\mathbb{N}}$, because every finite subset $J$ of
$\mathbb{N}$ satisfying $\left\{  2,3\right\}  \subseteq J\subseteq\mathbb{N}$
satisfies%
\[
\left[  x^{3}\right]  \left(  \sum_{i\in J}\left(  x+x^{2}\right)
^{i}\right)  =\left[  x^{3}\right]  \left(  \sum_{i\in\left\{  2,3\right\}
}\left(  x+x^{2}\right)  ^{i}\right)
\]
(this is simply a consequence of the fact that the only two entries of our
family that have a nonzero $x^{3}$-coefficient are the entries $\left(
x+x^{2}\right)  ^{i}$ for $i\in\left\{  2,3\right\}  $). Thus, any finite
subset of $\mathbb{N}$ that contains $\left\{  2,3\right\}  $ as a subset
determines the $x^{3}$-coefficient in the sum of this family $\left(  \left(
x+x^{2}\right)  ^{i}\right)  _{i\in\mathbb{N}}$. \medskip

\textbf{(b)} Consider the family%
\[
\left(  1+x^{i}\right)  _{i\in\mathbb{N}}=\left(  1+1,\ \ 1+x,\ \ 1+x^{2}%
,\ \ 1+x^{3},\ \ 1+x^{4},\ \ \ldots\right)
\]
of FPSs over $K=\mathbb{Z}$. The subset $\left\{  0,1,2,3\right\}  $ of
$\mathbb{N}$ determines the $x^{3}$-coefficient in the product of this family
$\left(  1+x^{i}\right)  _{i\in\mathbb{N}}$, because every finite subset $J$
of $\mathbb{N}$ satisfying $\left\{  0,1,2,3\right\}  \subseteq J\subseteq
\mathbb{N}$ satisfies%
\[
\left[  x^{3}\right]  \left(  \prod_{i\in J}\left(  1+x^{i}\right)  \right)
=\left[  x^{3}\right]  \left(  \prod_{i\in\left\{  0,1,2,3\right\}  }\left(
1+x^{i}\right)  \right)  .
\]
(This is because multiplying an FPS by any of the polynomials $1+x^{4}%
,\ \ 1+x^{5},\ \ 1+x^{6},\ \ \ldots$ leaves its $x^{3}$-coefficient
unchanged.) Thus, any finite subset of $\mathbb{N}$ that contains $\left\{
0,1,2,3\right\}  $ as a subset determines the $x^{3}$-coefficient in the
product of this family $\left(  1+x^{i}\right)  _{i\in\mathbb{N}}$.

On the other hand, the subset $\left\{  0,3\right\}  $ of $\mathbb{N}$ does
\textbf{not} determine the $x^{3}$-coefficient in the product of $\left(
1+x^{i}\right)  _{i\in\mathbb{N}}$. To see this, it suffices to notice that%
\[
\underbrace{\left[  x^{3}\right]  \left(  \prod_{i\in\left\{  0,1,2,3\right\}
}\left(  1+x^{i}\right)  \right)  }_{\substack{=\left[  x^{3}\right]  \left(
\left(  1+x^{0}\right)  \left(  1+x^{1}\right)  \left(  1+x^{2}\right)
\left(  1+x^{3}\right)  \right)  \\=4}}\neq\underbrace{\left[  x^{3}\right]
\left(  \prod_{i\in\left\{  0,3\right\}  }\left(  1+x^{i}\right)  \right)
}_{\substack{=\left[  x^{3}\right]  \left(  \left(  1+x^{0}\right)  \left(
1+x^{3}\right)  \right)  \\=2}}.
\]
(The philosophical reason is that, even though the monomial $x^{3}$ itself
does not appear in any of the entries $1+x^{1}$ and $1+x^{2}$, it does emerge
in the product of these two entries with the constant term of $\prod
_{i\in\left\{  0,3\right\}  }\left(  1+x^{i}\right)  =\left(  1+1\right)
\left(  1+x^{3}\right)  $.) \medskip

\textbf{(c)} Here is a simple but somewhat slippery example: Let $I=\left\{
1,2,3\right\}  $, and define the FPS $\mathbf{a}_{i}=1+\left(  -1\right)
^{i}x$ over $K=\mathbb{Z}$ for each $i\in I$ (so that $\mathbf{a}%
_{1}=\mathbf{a}_{3}=1-x$ and $\mathbf{a}_{2}=1+x$). Consider the finite family
$\left(  \mathbf{a}_{i}\right)  _{i\in I}$. Its product is already
well-defined by dint of its finiteness:%
\[
\prod_{i\in I}\mathbf{a}_{i}=\left(  1-x\right)  \left(  1+x\right)  \left(
1-x\right)  =1-x-x^{2}+x^{3}.
\]
But let us nevertheless see which subsets $M$ of $I$ determine the $x^{1}%
$-coefficient in the product of $\left(  \mathbf{a}_{i}\right)  _{i\in I}$.
The subset $I$ itself clearly does (since the only finite subset $J$ of $I$
satisfying $I\subseteq J\subseteq I$ is $I$ itself). The subset $\left\{
1\right\}  $ of $I$, however, does not, even though the product $\prod_{i\in
I}\mathbf{a}_{i}$ has the same $x^{1}$-coefficient as its subproduct
$\prod_{i\in\left\{  1\right\}  }\mathbf{a}_{i}$. (To see why $\left\{
1\right\}  $ does not determine the $x^{1}$-coefficient in the product of
$\left(  \mathbf{a}_{i}\right)  _{i\in I}$, it suffices to check that not
every finite subset $J$ of $I$ satisfying $\left\{  1\right\}  \subseteq
J\subseteq I$ satisfies the equality%
\[
\left[  x^{1}\right]  \left(  \prod_{i\in J}\mathbf{a}_{i}\right)  =\left[
x^{1}\right]  \left(  \prod_{i\in\left\{  1\right\}  }\mathbf{a}_{i}\right)
.
\]
For example, $J=\left\{  1,2\right\}  $ does not, since $\left[  x^{1}\right]
\left(  \prod_{i\in\left\{  1\right\}  }\mathbf{a}_{i}\right)  =-1$ but
$\left[  x^{1}\right]  \left(  \prod_{i\in\left\{  1,2\right\}  }%
\mathbf{a}_{i}\right)  =0$.)

This shows that, in order for a subset $M$ of $I$ to determine the $x^{n}%
$-coefficient in the product of a family $\left(  \mathbf{a}_{i}\right)
_{i\in I}$, it does \textbf{not} suffice to check that $\left[  x^{n}\right]
\left(  \prod_{i\in I}\mathbf{a}_{i}\right)  =\left[  x^{n}\right]  \left(
\prod_{i\in M}\mathbf{a}_{i}\right)  $; it rather needs to be shown that
$\left[  x^{n}\right]  \left(  \prod_{i\in J}\mathbf{a}_{i}\right)  =\left[
x^{n}\right]  \left(  \prod_{i\in M}\mathbf{a}_{i}\right)  $ for each finite
subset $J$ of $I$ satisfying $M\subseteq J\subseteq I$.
\end{example}

\begin{definition}
\label{def.fps.xn-coeff-fin-determined}Let $\left(  \mathbf{a}_{i}\right)
_{i\in I}\in K\left[  \left[  x\right]  \right]  ^{I}$ be a (possibly
infinite) family of FPSs. Let $n\in\mathbb{N}$. \medskip

\textbf{(a)} We say that \emph{the }$x^{n}$\emph{-coefficient in the sum of
}$\left(  \mathbf{a}_{i}\right)  _{i\in I}$\emph{ is finitely determined} if
there is a finite subset $M$ of $I$ that determines the $x^{n}$-coefficient in
the sum of $\left(  \mathbf{a}_{i}\right)  _{i\in I}$. \medskip

\textbf{(b)} We say that \emph{the }$x^{n}$\emph{-coefficient in the product
of }$\left(  \mathbf{a}_{i}\right)  _{i\in I}$\emph{ is finitely determined}
if there is a finite subset $M$ of $I$ that determines the $x^{n}$-coefficient
in the product of $\left(  \mathbf{a}_{i}\right)  _{i\in I}$.
\end{definition}

Using these concepts, we can now reword our definition of infinite sums as follows:

\begin{proposition}
\label{prop.fps.summable=fin-det}Let $\left(  \mathbf{a}_{i}\right)  _{i\in
I}\in K\left[  \left[  x\right]  \right]  ^{I}$ be a (possibly infinite)
family of FPSs. Then: \medskip

\textbf{(a)} The family $\left(  \mathbf{a}_{i}\right)  _{i\in I}$ is summable
if and only if each coefficient in its sum is finitely determined (i.e., for
each $n\in\mathbb{N}$, the $x^{n}$-coefficient in the sum of $\left(
\mathbf{a}_{i}\right)  _{i\in I}$ is finitely determined). \medskip

\textbf{(b)} If the family $\left(  \mathbf{a}_{i}\right)  _{i\in I}$ is
summable, then its sum $\sum_{i\in I}\mathbf{a}_{i}$ is the FPS whose $x^{n}%
$-coefficient (for any $n\in\mathbb{N}$) can be computed as follows: If
$n\in\mathbb{N}$, and if $M$ is a finite subset of $I$ that determines the
$x^{n}$-coefficient in the sum of $\left(  \mathbf{a}_{i}\right)  _{i\in I}$,
then
\[
\left[  x^{n}\right]  \left(  \sum_{i\in I}\mathbf{a}_{i}\right)  =\left[
x^{n}\right]  \left(  \sum_{i\in M}\mathbf{a}_{i}\right)  .
\]

\end{proposition}

\begin{proof}
Easy and LTTR.
\end{proof}

Inspired by Proposition \ref{prop.fps.summable=fin-det}, we can now define
infinite products of FPSs at last:

\begin{definition}
\label{def.fps.multipliable}Let $\left(  \mathbf{a}_{i}\right)  _{i\in I}$ be
a (possibly infinite) family of FPSs. Then: \medskip

\textbf{(a)} The family $\left(  \mathbf{a}_{i}\right)  _{i\in I}$ is said to
be \emph{multipliable} if and only if each coefficient in its product is
finitely determined. \medskip

\textbf{(b)} If the family $\left(  \mathbf{a}_{i}\right)  _{i\in I}$ is
multipliable, then its \emph{product} $\prod_{i\in I}\mathbf{a}_{i}$ is
defined to be the FPS whose $x^{n}$-coefficient (for any $n\in\mathbb{N}$) can
be computed as follows: If $n\in\mathbb{N}$, and if $M$ is a finite subset of
$I$ that determines the $x^{n}$-coefficient in the product of $\left(
\mathbf{a}_{i}\right)  _{i\in I}$, then
\[
\left[  x^{n}\right]  \left(  \prod_{i\in I}\mathbf{a}_{i}\right)  =\left[
x^{n}\right]  \left(  \prod_{i\in M}\mathbf{a}_{i}\right)  .
\]

\end{definition}

\begin{proposition}
\label{prop.fps.multipliable.prod-wd}This definition of $\prod_{i\in
I}\mathbf{a}_{i}$ is well-defined -- i.e., the coefficient $\left[
x^{n}\right]  \left(  \prod_{i\in M}\mathbf{a}_{i}\right)  $ does not depend
on $M$ (as long as $M$ is a finite subset of $I$ that determines the $x^{n}%
$-coefficient in the product of $\left(  \mathbf{a}_{i}\right)  _{i\in I}$).
\end{proposition}

\begin{proof}
Let $n\in\mathbb{N}$. We need to check that the coefficient $\left[
x^{n}\right]  \left(  \prod_{i\in M}\mathbf{a}_{i}\right)  $ does not depend
on $M$ (as long as $M$ is a finite subset of $I$ that determines the $x^{n}%
$-coefficient in the product of $\left(  \mathbf{a}_{i}\right)  _{i\in I}$).
In other words, we need to check that if $M_{1}$ and $M_{2}$ are two finite
subsets of $I$ that each determine the $x^{n}$-coefficient in the product of
$\left(  \mathbf{a}_{i}\right)  _{i\in I}$, then%
\begin{equation}
\left[  x^{n}\right]  \left(  \prod_{i\in M_{1}}\mathbf{a}_{i}\right)
=\left[  x^{n}\right]  \left(  \prod_{i\in M_{2}}\mathbf{a}_{i}\right)  .
\label{pf.prop.fps.multipliable.prod-wd.goal}%
\end{equation}


So let us prove this. Let $M_{1}$ and $M_{2}$ be two finite subsets of $I$
that each determine the $x^{n}$-coefficient in the product of $\left(
\mathbf{a}_{i}\right)  _{i\in I}$. Thus, in particular, $M_{1}$ determines the
$x^{n}$-coefficient in the product of $\left(  \mathbf{a}_{i}\right)  _{i\in
I}$. In other words, every finite subset $J$ of $I$ satisfying $M_{1}\subseteq
J\subseteq I$ satisfies%
\[
\left[  x^{n}\right]  \left(  \prod_{i\in J}\mathbf{a}_{i}\right)  =\left[
x^{n}\right]  \left(  \prod_{i\in M_{1}}\mathbf{a}_{i}\right)  .
\]
Applying this to $J=M_{1}\cup M_{2}$, we obtain%
\begin{equation}
\left[  x^{n}\right]  \left(  \prod_{i\in M_{1}\cup M_{2}}\mathbf{a}%
_{i}\right)  =\left[  x^{n}\right]  \left(  \prod_{i\in M_{1}}\mathbf{a}%
_{i}\right)  \label{pf.prop.fps.multipliable.prod-wd.1}%
\end{equation}
(since $M_{1}\cup M_{2}$ is a subset of $I$ satisfying $M_{1}\subseteq
M_{1}\cup M_{2}\subseteq I$). The same argument (with the roles of $M_{1}$ and
$M_{2}$ swapped) yields%
\begin{equation}
\left[  x^{n}\right]  \left(  \prod_{i\in M_{2}\cup M_{1}}\mathbf{a}%
_{i}\right)  =\left[  x^{n}\right]  \left(  \prod_{i\in M_{2}}\mathbf{a}%
_{i}\right)  . \label{pf.prop.fps.multipliable.prod-wd.2}%
\end{equation}
The left hand sides of the equalities
(\ref{pf.prop.fps.multipliable.prod-wd.1}) and
(\ref{pf.prop.fps.multipliable.prod-wd.2}) are equal (since $M_{1}\cup
M_{2}=M_{2}\cup M_{1}$). Thus, the right hand sides are equal as well. In
other words, $\left[  x^{n}\right]  \left(  \prod_{i\in M_{1}}\mathbf{a}%
_{i}\right)  =\left[  x^{n}\right]  \left(  \prod_{i\in M_{2}}\mathbf{a}%
_{i}\right)  $. Thus, we have proved
(\ref{pf.prop.fps.multipliable.prod-wd.goal}), and with it Proposition
\ref{prop.fps.multipliable.prod-wd}.
\end{proof}

The attentive (and pedantic) reader might notice that there is one more thing
that needs to be checked in order to make sure that Definition
\ref{def.fps.multipliable} \textbf{(b)} is legitimate. In fact, this
definition does not merely define (some) infinite products $\prod_{i\in
I}\mathbf{a}_{i}$ of FPSs, but also \textquotedblleft
accidentally\textquotedblright\ gives a new meaning to \textbf{finite}
products $\prod_{i\in I}\mathbf{a}_{i}$ (since a finite family $\left(
\mathbf{a}_{i}\right)  _{i\in I}$ of FPSs is always multipliable). We
therefore need to check that this new meaning does not conflict with the
original definition of a finite product of elements of a commutative ring. In
other words, we need to prove the following:

\begin{proposition}
\label{prop.fps.multipliable.prod-wd2}Let $\left(  \mathbf{a}_{i}\right)
_{i\in I}$ be a finite family of FPSs. Then, the product $\prod_{i\in
I}\mathbf{a}_{i}$ defined according to Definition \ref{def.fps.multipliable}
\textbf{(b)} equals the finite product $\prod_{i\in I}\mathbf{a}_{i}$ defined
in the usual way (i.e., defined as in any commutative ring).
\end{proposition}

\begin{proof}
Argue that $I$ itself is a subset of $I$ that determines all coefficients in
the product of $\left(  \mathbf{a}_{i}\right)  _{i\in I}$. See Section
\ref{sec.details.gf.prod} for a detailed proof.
\end{proof}

We shall apply Convention \ref{conv.fps.infsum} to infinite products just like
we have been applying it to infinite sums. For instance, the product sign
$\prod\limits_{k=m}^{\infty}$ (for a fixed $m\in\mathbb{Z}$) means
$\prod\limits_{k\in\left\{  m,m+1,m+2,\ldots\right\}  }$.

\subsubsection{Why $\prod_{i\in\mathbb{N}}\left(  1+x^{2^{i}}\right)  $ works
and $\prod_{i\in\mathbb{N}}\left(  1+ix\right)  $ doesn't}

Let us now see how Definition \ref{def.fps.multipliable} legitimizes our
product $\prod_{i\in\mathbb{N}}\left(  1+x^{2^{i}}\right)  $ from Subsection
\ref{subsec.gf.prod.exa}. Indeed,%
\[
\prod_{i\in\mathbb{N}}\left(  1+x^{2^{i}}\right)  =\left(  1+x^{1}\right)
\left(  1+x^{2}\right)  \left(  1+x^{4}\right)  \left(  1+x^{8}\right)
\cdots.
\]
If you want to compute the $x^{6}$-coefficient in this product, you only need
to multiply the first 3 factors $\left(  1+x^{1}\right)  \left(
1+x^{2}\right)  \left(  1+x^{4}\right)  $; none of the other factors will
change this coefficient in any way, because multiplying an FPS by $1+x^{m}$
(for some $m>0$) does not change its first $m$ coefficients\footnote{For
example, let us check this for $m=3$: If we multiply an FPS $a_{0}x^{0}%
+a_{1}x^{1}+a_{2}x^{2}+\cdots$ by $1+x^{3}$, then we obtain%
\begin{align*}
&  \left(  a_{0}x^{0}+a_{1}x^{1}+a_{2}x^{2}+\cdots\right)  \left(
1+x^{3}\right) \\
&  =a_{0}x^{0}+a_{1}x^{1}+a_{2}x^{2}+\left(  a_{3}+a_{0}\right)  x^{3}+\left(
a_{4}+a_{1}\right)  x^{4}+\left(  a_{5}+a_{2}\right)  x^{5}+\cdots,
\end{align*}
and so the first $3$ coefficients are left unchanged.}. Likewise, if you want
to compute the $x^{13}$-coefficient of the above product, then you only need
to multiply the first $4$ factors; none of the others will have any effect on
this coefficient. The same logic applies to the $x^{n}$-coefficient for any
$n\in\mathbb{N}$; it is determined by the first $\left\lfloor \log
_{2}n\right\rfloor +1$ factors of the product. Thus, each coefficient in the
product is finitely determined. This means that the family is multipliable;
thus, its product makes sense.

In contrast, the product%
\[
\left(  1+0x\right)  \left(  1+1x\right)  \left(  1+2x\right)  \left(
1+3x\right)  \left(  1+4x\right)  \cdots=\prod_{i\in\mathbb{N}}\left(
1+ix\right)
\]
does not make sense. Indeed, its $x^{1}$-coefficient is not finitely
determined (any of the factors other than $1+0x$ affects it), so the family
$\left(  1+ix\right)  _{i\in\mathbb{N}}$ is not multipliable.

Likewise, the product%
\[
\left(  1+\dfrac{x}{1^{2}}\right)  \left(  1+\dfrac{x}{2^{2}}\right)  \left(
1+\dfrac{x}{3^{2}}\right)  \cdots=\prod_{i\in\left\{  1,2,3,\ldots\right\}
}\left(  1+\dfrac{x}{i^{2}}\right)
\]
does not make sense in $K\left[  \left[  x\right]  \right]  $ (although the
analogous product in complex analysis defines a holomorphic function).

\subsubsection{A general criterion for multipliability}

Recall our reasoning that we used above to prove that the family $\left(
1+x^{2^{i}}\right)  _{i\in\mathbb{N}}$ is multipliable. The core of this
reasoning was the observation that multiplying an FPS by $1+x^{m}$ (for some
$m>0$) does not change its first $m$ coefficients. This can be generalized: If
$f\in K\left[  \left[  x\right]  \right]  $ is an FPS whose first $m$
coefficients are $0$ (for example, $f$ can be $x^{m}$, in which case we
recover the statement in our preceding sentence), then multiplying an FPS $a$
by $1+f$ does not change its first $m$ coefficients (that is, the first $m$
coefficients of $a\left(  1+f\right)  $ are the first $m$ coefficients of
$a$). This is a useful fact, so let us state it as a lemma (renaming $m$ as
$n+1$):

\begin{lemma}
\label{lem.fps.prod.irlv.1}Let $a,f\in K\left[  \left[  x\right]  \right]  $
be two FPSs. Let $n\in\mathbb{N}$. Assume that
\begin{equation}
\left[  x^{m}\right]  f=0\ \ \ \ \ \ \ \ \ \ \text{for each }m\in\left\{
0,1,\ldots,n\right\}  . \label{eq.lem.fps.prod.irlv.1.ass}%
\end{equation}
Then,
\[
\left[  x^{m}\right]  \left(  a\left(  1+f\right)  \right)  =\left[
x^{m}\right]  a\ \ \ \ \ \ \ \ \ \ \text{for each }m\in\left\{  0,1,\ldots
,n\right\}  .
\]

\end{lemma}

\begin{proof}
[Proof of Lemma \ref{lem.fps.prod.irlv.1}.]The FPS $af$ is a multiple of $f$
(since $af=fa$). Hence, Lemma \ref{lem.fps.prod.irlv.mul} (applied to $u=f$
and $v=af$) yields that%
\begin{equation}
\left[  x^{m}\right]  \left(  af\right)  =0\ \ \ \ \ \ \ \ \ \ \text{for each
}m\in\left\{  0,1,\ldots,n\right\}  \label{pf.lem.fps.prod.irlv.1.4}%
\end{equation}
(since we have assumed that $\left[  x^{m}\right]  f=0$ for each $m\in\left\{
0,1,\ldots,n\right\}  $).

Now, for each $m\in\left\{  0,1,\ldots,n\right\}  $, we have
\begin{align*}
\left[  x^{m}\right]  \left(  \underbrace{a\left(  1+f\right)  }%
_{=a+af}\right)   &  =\left[  x^{m}\right]  \left(  a+af\right)  =\left[
x^{m}\right]  a+\underbrace{\left[  x^{m}\right]  \left(  af\right)
}_{\substack{=0\\\text{(by (\ref{pf.lem.fps.prod.irlv.1.4}))}}%
}\ \ \ \ \ \ \ \ \ \ \left(  \text{by (\ref{pf.thm.fps.ring.xn(a+b)=})}\right)
\\
&  =\left[  x^{m}\right]  a.
\end{align*}
This proves Lemma \ref{lem.fps.prod.irlv.1}.
\end{proof}

For convenience, let us extend Lemma \ref{lem.fps.prod.irlv.1} to products of
several factors:

\begin{lemma}
\label{lem.fps.prod.irlv.fin}Let $a\in K\left[  \left[  x\right]  \right]  $
be an FPS. Let $\left(  f_{i}\right)  _{i\in J}\in K\left[  \left[  x\right]
\right]  ^{J}$ be a finite family of FPSs. Let $n\in\mathbb{N}$. Assume that
each $i\in J$ satisfies
\begin{equation}
\left[  x^{m}\right]  \left(  f_{i}\right)  =0\ \ \ \ \ \ \ \ \ \ \text{for
each }m\in\left\{  0,1,\ldots,n\right\}  .
\label{eq.lem.fps.prod.irlv.fin.ass}%
\end{equation}
Then,
\[
\left[  x^{m}\right]  \left(  a\prod_{i\in J}\left(  1+f_{i}\right)  \right)
=\left[  x^{m}\right]  a\ \ \ \ \ \ \ \ \ \ \text{for each }m\in\left\{
0,1,\ldots,n\right\}  .
\]

\end{lemma}

\begin{proof}
[Proof of Lemma \ref{lem.fps.prod.irlv.fin}.]This is just Lemma
\ref{lem.fps.prod.irlv.1}, applied several times (specifically, $\left\vert
J\right\vert $ many times). See Section \ref{sec.details.gf.prod} for a
detailed proof.
\end{proof}

Now, using Lemma \ref{lem.fps.prod.irlv.fin}, we can obtain the following
convenient criterion for multipliability:

\begin{theorem}
\label{thm.fps.1+f-mulable}Let $\left(  f_{i}\right)  _{i\in I}\in K\left[
\left[  x\right]  \right]  ^{I}$ be a (possibly infinite) summable family of
FPSs. Then, the family $\left(  1+f_{i}\right)  _{i\in I}$ is multipliable.
\end{theorem}

\begin{proof}
[Proof of Theorem \ref{thm.fps.1+f-mulable}.]This is an easy consequence of
Lemma \ref{lem.fps.prod.irlv.fin}. See Section \ref{sec.details.gf.prod} for a
detailed proof.
\end{proof}

We notice two simple sufficient (if rarely satisfied) criteria for multipliability:

\begin{proposition}
\label{prop.fps.1-mulable}If all but finitely many entries of a family
$\left(  \mathbf{a}_{i}\right)  _{i\in I}\in K\left[  \left[  x\right]
\right]  ^{I}$ equal $1$ (that is, if all but finitely many $i\in I$ satisfy
$\mathbf{a}_{i}=1$), then this family is multipliable.
\end{proposition}

\begin{proof}
LTTR. (See Section \ref{sec.details.gf.prod} for a detailed proof.)
\end{proof}

\begin{remark}
\label{rmk.fps.0-mulable}If a family $\left(  \mathbf{a}_{i}\right)  _{i\in
I}\in K\left[  \left[  x\right]  \right]  ^{I}$ contains $0$ as an entry
(i.e., if there exists an $i\in I$ such that $\mathbf{a}_{i}=0$), then this
family is automatically multipliable, and its product is $0$.
\end{remark}

\begin{proof}
Assume that the family $\left(  \mathbf{a}_{i}\right)  _{i\in I}$ contains $0$
as an entry. That is, there exists some $j\in I$ such that $\mathbf{a}_{j}=0$.
Consider this $j$. Now, it is easy to see that the subset $\left\{  j\right\}
$ of $I$ determines all coefficients in the product of $\left(  \mathbf{a}%
_{i}\right)  _{i\in I}$. The details are LTTR.
\end{proof}

\subsubsection{$x^{n}$-approximators}

Working with multipliable families gets slightly easier using the following notion:

\begin{definition}
\label{def.fps.infprod-approx}Let $\left(  \mathbf{a}_{i}\right)  _{i\in I}\in
K\left[  \left[  x\right]  \right]  ^{I}$ be a family of FPSs. Let
$n\in\mathbb{N}$. An $x^{n}$\emph{-approximator} for $\left(  \mathbf{a}%
_{i}\right)  _{i\in I}$ means a finite subset $M$ of $I$ that determines the
first $n+1$ coefficients in the product of $\left(  \mathbf{a}_{i}\right)
_{i\in I}$. (In other words, $M$ has to determine the $x^{m}$-coefficient in
the product of $\left(  \mathbf{a}_{i}\right)  _{i\in I}$ for each
$m\in\left\{  0,1,\ldots,n\right\}  $.)
\end{definition}

The name \textquotedblleft$x^{n}$-approximator\textquotedblright\ is supposed
to hint at the fact that if $M$ is an $x^{n}$-approximator for a multipliable
family $\left(  \mathbf{a}_{i}\right)  _{i\in I}$, then the (finite)
subproduct $\prod_{i\in M}\mathbf{a}_{i}$ \textquotedblleft
approximates\textquotedblright\ the full product $\prod_{i\in I}\mathbf{a}%
_{i}$ up until the $x^{n}$-coefficient (i.e., the first $n+1$ coefficients of
$\prod_{i\in M}\mathbf{a}_{i}$ equal the respective coefficients of
$\prod_{i\in I}\mathbf{a}_{i}$). See Proposition
\ref{prop.fps.infprod-approx-xneq} \textbf{(b)} below for the precise
statement of this fact.

Clearly, an $x^{n}$-approximator for a family $\left(  \mathbf{a}_{i}\right)
_{i\in I}$ always determines the $x^{n}$-coefficient in the product of
$\left(  \mathbf{a}_{i}\right)  _{i\in I}$. But the converse is not true, as
the following example shows:

\begin{example}
Consider the family
\[
\left(  1+x^{2^{i}}\right)  _{i\in\mathbb{N}}=\left(  1+x^{1},\ \ 1+x^{2}%
,\ \ 1+x^{4},\ \ 1+x^{8},\ \ \ldots\right)
\]
of FPSs. The finite subset $\left\{  1,2\right\}  $ of $\mathbb{N}$ determines
the $x^{6}$-coefficient in the product of this family (indeed, the $x^{6}%
$-coefficient of the product $\left(  1+x^{2}\right)  \left(  1+x^{4}\right)
$ is $1$, and this does not change if we multiply any further factors onto
this product), but is not an $x^{6}$-approximator for this family (since,
e.g., it does not determine the $x^{5}$-coefficient in its product).
\end{example}

\begin{lemma}
\label{lem.fps.mulable.approx}Let $\left(  \mathbf{a}_{i}\right)  _{i\in I}\in
K\left[  \left[  x\right]  \right]  ^{I}$ be a multipliable family of FPSs.
Let $n\in\mathbb{N}$. Then, there exists an $x^{n}$-approximator for $\left(
\mathbf{a}_{i}\right)  _{i\in I}$.
\end{lemma}

\begin{proof}
[Proof of Lemma \ref{lem.fps.mulable.approx} (sketched).]This is an easy
consequence of the fact that a union of finitely many finite sets is finite. A
detailed proof can be found in Section \ref{sec.details.gf.prod}.
\end{proof}

As promised above, we can use $x^{n}$-approximators to \textquotedblleft
approximate\textquotedblright\ infinite products of FPSs (in the sense of:
compute the first $n+1$ coefficients of these products). Here is why this
works:\footnote{See Definition \ref{def.fps.xneq} for the meaning of the
symbol \textquotedblleft$\overset{x^{n}}{\equiv}$\textquotedblright\ appearing
in this proposition.}

\begin{proposition}
\label{prop.fps.infprod-approx-xneq}Let $\left(  \mathbf{a}_{i}\right)  _{i\in
I}\in K\left[  \left[  x\right]  \right]  ^{I}$ be a family of FPSs. Let
$n\in\mathbb{N}$. Let $M$ be an $x^{n}$-approximator for $\left(
\mathbf{a}_{i}\right)  _{i\in I}$. Then: \medskip

\textbf{(a)} Every finite subset $J$ of $I$ satisfying $M\subseteq J\subseteq
I$ satisfies%
\[
\prod_{i\in J}\mathbf{a}_{i}\overset{x^{n}}{\equiv}\prod_{i\in M}%
\mathbf{a}_{i}.
\]


\textbf{(b)} If the family $\left(  \mathbf{a}_{i}\right)  _{i\in I}$ is
multipliable, then
\[
\prod_{i\in I}\mathbf{a}_{i}\overset{x^{n}}{\equiv}\prod_{i\in M}%
\mathbf{a}_{i}.
\]

\end{proposition}

\begin{proof}
This follows easily from Definition \ref{def.fps.infprod-approx} and
Definition \ref{def.fps.multipliable} \textbf{(b)}. See Section
\ref{sec.details.gf.prod} for a detailed proof.
\end{proof}

\subsubsection{Properties of infinite products}

We shall now establish some general properties of infinite products. By and
large, these properties are analogous to corresponding properties of finite
products, although they need a few technical requirements (see Remark
\ref{rmk.fps.subfamily-not-mulable} for why).

We begin with the first property, which says in essence that a product can be
broken into two parts:

\begin{proposition}
\label{prop.fps.union-mulable}Let $\left(  \mathbf{a}_{i}\right)  _{i\in I}\in
K\left[  \left[  x\right]  \right]  ^{I}$ be a family of FPSs. Let $J$ be a
subset of $I$. Assume that the subfamilies $\left(  \mathbf{a}_{i}\right)
_{i\in J}$ and $\left(  \mathbf{a}_{i}\right)  _{i\in I\setminus J}$ are
multipliable. Then: \medskip

\textbf{(a)} The entire family $\left(  \mathbf{a}_{i}\right)  _{i\in I}$ is
multipliable. \medskip

\textbf{(b)} We have
\[
\prod_{i\in I}\mathbf{a}_{i}=\left(  \prod_{i\in J}\mathbf{a}_{i}\right)
\cdot\left(  \prod_{i\in I\setminus J}\mathbf{a}_{i}\right)  .
\]

\end{proposition}

\begin{proof}
[Proof of Proposition \ref{prop.fps.union-mulable} (sketched).]Here is the
idea: Fix $n\in\mathbb{N}$. Lemma \ref{lem.fps.mulable.approx} (applied to $J$
instead of $I$) shows that there exists an $x^{n}$-approximator $U$ for
$\left(  \mathbf{a}_{i}\right)  _{i\in J}$. Consider this $U$. Lemma
\ref{lem.fps.mulable.approx} (applied to $J$ instead of $I$) shows that there
exists an $x^{n}$-approximator $V$ for $\left(  \mathbf{a}_{i}\right)  _{i\in
I\setminus J}$. Consider this $V$. Note that $U\cup V$ is finite (since $U$
and $V$ are finite). Now, it is not hard to see that $U\cup V$ determines the
$x^{n}$-coefficient in the product of $\left(  \mathbf{a}_{i}\right)  _{i\in
I}$ (indeed, it is not much harder to see that $U\cup V$ is an $x^{n}%
$-approximator for $\left(  \mathbf{a}_{i}\right)  _{i\in I}$). Hence, the
$x^{n}$-coefficient in the product of $\left(  \mathbf{a}_{i}\right)  _{i\in
I}$ is finitely determined (since $U\cup V$ is finite). Now, forget that we
fixed $n$, and conclude that the family $\left(  \mathbf{a}_{i}\right)  _{i\in
I}$ is multipliable. This proves part \textbf{(a)}. Part \textbf{(b)} easily
follows using Proposition \ref{prop.fps.infprod-approx-xneq} \textbf{(b)}.

The details of this proof can be found in Section \ref{sec.details.gf.prod}.
\end{proof}

A useful particular case of Proposition \ref{prop.fps.union-mulable} is
obtained when the subset $J$ is a single-element set $\left\{  j\right\}  $.
In this case, the proposition says that%
\[
\prod_{i\in I}\mathbf{a}_{i}=\mathbf{a}_{j}\cdot\prod_{i\in I\setminus\left\{
j\right\}  }\mathbf{a}_{i}%
\]
(assuming that the family $\left(  \mathbf{a}_{i}\right)  _{i\in
I\setminus\left\{  j\right\}  }$ is multipliable\footnote{The one-element
family $\left(  \mathbf{a}_{i}\right)  _{i\in\left\{  j\right\}  }$ is, of
course, always multipliable.}). This rule allows us to split off any factor
from a multipliable product, as long as the rest of the product is still
multipliable. \medskip

Our next property generalizes the classical rule $\prod_{i\in I}\left(
a_{i}b_{i}\right)  =\left(  \prod_{i\in I}a_{i}\right)  \cdot\left(
\prod_{i\in I}b_{i}\right)  $ of finite products to infinite ones:

\begin{proposition}
\label{prop.fps.prod-mulable}Let $\left(  \mathbf{a}_{i}\right)  _{i\in I}\in
K\left[  \left[  x\right]  \right]  ^{I}$ and $\left(  \mathbf{b}_{i}\right)
_{i\in I}\in K\left[  \left[  x\right]  \right]  ^{I}$ be two multipliable
families of FPSs. Then: \medskip

\textbf{(a)} The family $\left(  \mathbf{a}_{i}\mathbf{b}_{i}\right)  _{i\in
I}$ is multipliable. \medskip

\textbf{(b)} We have%
\[
\prod_{i\in I}\left(  \mathbf{a}_{i}\mathbf{b}_{i}\right)  =\left(
\prod_{i\in I}\mathbf{a}_{i}\right)  \cdot\left(  \prod_{i\in I}\mathbf{b}%
_{i}\right)  .
\]

\end{proposition}

\begin{proof}
[Proof of Proposition \ref{prop.fps.prod-mulable} (sketched).]Here is the
idea: Fix $n\in\mathbb{N}$. Lemma \ref{lem.fps.mulable.approx} shows that
there exists an $x^{n}$-approximator $U$ for $\left(  \mathbf{a}_{i}\right)
_{i\in I}$. Consider this $U$. Lemma \ref{lem.fps.mulable.approx} (applied to
$\mathbf{b}_{i}$ instead of $\mathbf{a}_{i}$) shows that there exists an
$x^{n}$-approximator $V$ for $\left(  \mathbf{b}_{i}\right)  _{i\in I}$.
Consider this $V$. Note that $U\cup V$ is finite (since $U$ and $V$ are
finite). Now, it is not hard to see that $U\cup V$ is an $x^{n}$-approximator
for $\left(  \mathbf{a}_{i}\mathbf{b}_{i}\right)  _{i\in I}$. From here,
proceed as in the proof of Proposition \ref{prop.fps.union-mulable}.

The details of this proof can be found in Section \ref{sec.details.gf.prod}.
\end{proof}

Next comes an analogous property for quotients instead of products:

\begin{proposition}
\label{prop.fps.div-mulable}Let $\left(  \mathbf{a}_{i}\right)  _{i\in I}\in
K\left[  \left[  x\right]  \right]  ^{I}$ and $\left(  \mathbf{b}_{i}\right)
_{i\in I}\in K\left[  \left[  x\right]  \right]  ^{I}$ be two multipliable
families of FPSs. Assume that the FPS $\mathbf{b}_{i}$ is invertible for each
$i\in I$. Then: \medskip

\textbf{(a)} The family $\left(  \dfrac{\mathbf{a}_{i}}{\mathbf{b}_{i}%
}\right)  _{i\in I}$ is multipliable. \medskip

\textbf{(b)} We have%
\[
\prod_{i\in I}\dfrac{\mathbf{a}_{i}}{\mathbf{b}_{i}}=\dfrac{\prod\limits_{i\in
I}\mathbf{a}_{i}}{\prod\limits_{i\in I}\mathbf{b}_{i}}.
\]

\end{proposition}

\begin{proof}
[Proof of Proposition \ref{prop.fps.div-mulable} (sketch).]This is similar to
the proof of Proposition \ref{prop.fps.prod-mulable}, but using
(\ref{eq.thm.fps.xneq.props.e./}) instead of Lemma
\ref{lem.fps.prod.irlv.cong-mul}. The details of this proof can be found in
Section \ref{sec.details.gf.prod}.
\end{proof}

Now we come to an annoying technicality. In Proposition
\ref{prop.fps.summable.sub}, we have learnt that any subfamily of a summable
family is again summable. Alas, the analogous fact for multipliability is false:

\begin{remark}
\label{rmk.fps.subfamily-not-mulable}Not every subfamily of a multipliable
family is itself multipliable. For example, for $K=\mathbb{Z}$, the family
$\left(  0,1,2,3,\ldots\right)  $ is multipliable, but its subfamily $\left(
1,2,3,\ldots\right)  $ is not.
\end{remark}

Nevertheless, if a family of \textbf{invertible} FPSs is multipliable, then so
is any subfamily of it. In other words:

\begin{proposition}
\label{prop.fps.prods-mulable-subfams}Let $\left(  \mathbf{a}_{i}\right)
_{i\in I}\in K\left[  \left[  x\right]  \right]  ^{I}$ be a multipliable
family of invertible FPSs. Then, any subfamily of $\left(  \mathbf{a}%
_{i}\right)  _{i\in I}$ is multipliable.
\end{proposition}

\begin{proof}
[Proof of Proposition \ref{prop.fps.prods-mulable-subfams} (sketched).]This is
another proof in the tradition of the proofs of Proposition
\ref{prop.fps.union-mulable} and Proposition \ref{prop.fps.prod-mulable}. We
must show that the family $\left(  \mathbf{a}_{i}\right)  _{i\in J}$ is
multipliable whenever $J$ is a subset of $I$. The idea is to show that if $U$
is an $x^{n}$-approximator for $\left(  \mathbf{a}_{i}\right)  _{i\in I}$,
then $U\cap J$ determines the $x^{n}$-coefficient in the product of $\left(
\mathbf{a}_{i}\right)  _{i\in J}$ (and, in fact, is an $x^{n}$-approximator
for $\left(  \mathbf{a}_{i}\right)  _{i\in J}$). This relies on the
invertibility of $\prod_{i\in U\setminus J}\mathbf{a}_{i}$; this is why the
FPS $\mathbf{a}_{i}$ are required to be invertible in the proposition.

The details of this proof can be found in Section \ref{sec.details.gf.prod}.
\end{proof}

As Remark \ref{rmk.fps.subfamily-not-mulable} shows, Proposition
\ref{prop.fps.prods-mulable-subfams} would not hold without the word
\textquotedblleft invertible\textquotedblright. \medskip

Now let us state a trivial rule that is nevertheless worth stating. Recall
that finite products can be reindexed using a bijection -- i.e., if
$f:S\rightarrow T$ is a bijection between two finite sets $S$ and $T$, then
any product $\prod_{t\in T}\mathbf{a}_{t}$ can be rewritten as $\prod_{s\in
S}\mathbf{a}_{f\left(  s\right)  }$. The same holds for infinite products:

\begin{proposition}
\label{prop.fps.prods-mulable-rules.reindex}Let $S$ and $T$ be two sets. Let
$f:S\rightarrow T$ be a bijection. Let $\left(  \mathbf{a}_{t}\right)  _{t\in
T}\in K\left[  \left[  x\right]  \right]  ^{T}$ be a multipliable family of
FPSs. Then,
\[
\prod_{t\in T}\mathbf{a}_{t}=\prod_{s\in S}\mathbf{a}_{f\left(  s\right)  }%
\]
(and, in particular, the product on the right hand side is well-defined, i.e.,
the family $\left(  \mathbf{a}_{f\left(  s\right)  }\right)  _{s\in S}$ is multipliable).
\end{proposition}

In other words, if we reindex a multipliable family of FPSs (using a
bijection), then the resulting family will still be multipliable and have the
same product as the original family.

\begin{proof}
[Proof of Proposition \ref{prop.fps.prods-mulable-rules.reindex}
(sketched).]This is a trivial consequence of the definitions of
multipliability and infinite products.
\end{proof}

Next comes a rule that allows us to break an infinite product into any amount
(possibly infinite!) of subproducts:

\begin{proposition}
\label{prop.fps.prods-mulable-rules.SW1}Let $\left(  \mathbf{a}_{s}\right)
_{s\in S}\in K\left[  \left[  x\right]  \right]  ^{S}$ be a multipliable
family of FPSs. Let $W$ be a set. Let $f:S\rightarrow W$ be a map. Assume that
for each $w\in W$, the family $\left(  \mathbf{a}_{s}\right)  _{s\in
S;\ f\left(  s\right)  =w}$ is multipliable.\footnotemark\ Then,%
\begin{equation}
\prod_{s\in S}\mathbf{a}_{s}=\prod_{w\in W}\ \ \prod_{\substack{s\in
S;\\f\left(  s\right)  =w}}\mathbf{a}_{s}.
\label{eq.prop.fps.prods-mulable-rules.SW1.eq}%
\end{equation}
(In particular, the right hand side is well-defined -- i.e., the family
$\left(  \prod_{\substack{s\in S;\\f\left(  s\right)  =w}}\mathbf{a}%
_{s}\right)  _{w\in W}$ is multipliable.)
\end{proposition}

\footnotetext{Note that this assumption automatically holds if we assume that
all FPSs $\mathbf{a}_{s}$ are invertible. Indeed, the family $\left(
\mathbf{a}_{s}\right)  _{s\in S;\ f\left(  s\right)  =w}$ is a subfamily of
the multipliable family $\left(  \mathbf{a}_{s}\right)  _{s\in S}$ and thus
must itself be multipliable if all the $\mathbf{a}_{s}$ are invertible (by
Proposition \ref{prop.fps.prods-mulable-subfams}).}

\begin{proof}
[Proof of Proposition \ref{prop.fps.prods-mulable-rules.SW1} (sketched).]This
can be derived from the analogous property of finite products, since all
coefficients in a multipliable product are finitely determined (conveniently
using $x^{n}$-approximators, which determine several coefficients at the same
time). See Section \ref{sec.details.gf.prod} for the details of this proof.
\end{proof}

The next rule allows (under certain technical conditions) to interchange
product signs, and to rewrite a nested product as a product over pairs:

\begin{proposition}
[Fubini rule for infinite products of FPSs]%
\label{prop.fps.prods-mulable-rules.fubini1}Let $I$ and $J$ be two sets. Let
$\left(  \mathbf{a}_{\left(  i,j\right)  }\right)  _{\left(  i,j\right)  \in
I\times J}\in K\left[  \left[  x\right]  \right]  ^{I\times J}$ be a
multipliable family of FPSs. Assume that for each $i\in I$, the family
$\left(  \mathbf{a}_{\left(  i,j\right)  }\right)  _{j\in J}$ is multipliable.
Assume that for each $j\in J$, the family $\left(  \mathbf{a}_{\left(
i,j\right)  }\right)  _{i\in I}$ is multipliable.\ Then,%
\[
\prod_{i\in I}\ \ \prod_{j\in J}\mathbf{a}_{\left(  i,j\right)  }%
=\prod_{\left(  i,j\right)  \in I\times J}\mathbf{a}_{\left(  i,j\right)
}=\prod_{j\in J}\mathbf{\ \ }\prod_{i\in I}\mathbf{a}_{\left(  i,j\right)  }.
\]
(In particular, all the products appearing in this equality are well-defined.)
\end{proposition}

\begin{proof}
[Proof of Proposition \ref{prop.fps.prods-mulable-rules.fubini1}
(sketched).]The first equality follows by applying Proposition
\ref{prop.fps.prods-mulable-rules.SW1} to $S=I\times J$ and $W=I$ and
$f\left(  i,j\right)  =i$ (and appropriately reindexing the products). The
second equality is analogous. See Section \ref{sec.details.gf.prod} for the
details of this proof.
\end{proof}

The above rules (Proposition \ref{prop.fps.union-mulable}, Proposition
\ref{prop.fps.prod-mulable}, Proposition \ref{prop.fps.div-mulable},
Proposition \ref{prop.fps.prods-mulable-subfams}, Proposition
\ref{prop.fps.prods-mulable-rules.reindex}, Proposition
\ref{prop.fps.prods-mulable-rules.SW1}, Proposition
\ref{prop.fps.prods-mulable-rules.fubini1}) show that infinite products (as we
have defined them) are well-behaved -- i.e., satisfy the usual rules that
finite products satisfy, with only one minor caveat: Subfamilies of
multipliable families might fail to be multipliable (as saw in Remark
\ref{rmk.fps.subfamily-not-mulable}). If we restrict ourselves to multipliable
families of \textbf{invertible} FPSs, then even this caveat is avoided:

\begin{proposition}
[Fubini rule for infinite products of FPSs, invertible case]%
\label{prop.fps.prods-mulable-rules.fubini}Let $I$ and $J$ be two sets. Let
$\left(  \mathbf{a}_{\left(  i,j\right)  }\right)  _{\left(  i,j\right)  \in
I\times J}\in K\left[  \left[  x\right]  \right]  ^{I\times J}$ be a
multipliable family of invertible FPSs. Then,%
\[
\prod_{i\in I}\ \ \prod_{j\in J}\mathbf{a}_{\left(  i,j\right)  }%
=\prod_{\left(  i,j\right)  \in I\times J}\mathbf{a}_{\left(  i,j\right)
}=\prod_{j\in J}\mathbf{\ \ }\prod_{i\in I}\mathbf{a}_{\left(  i,j\right)  }.
\]
(In particular, all the products appearing in this equality are well-defined.)
\end{proposition}

\begin{proof}
[Proof of Proposition \ref{prop.fps.prods-mulable-rules.fubini} (sketched).]%
See Section \ref{sec.details.gf.prod}.
\end{proof}

These rules (and some similar ones, which the reader can easily invent and
prove\footnote{For example, if $\left(  \mathbf{a}_{i}\right)  _{i\in I}$ is a
multipliable family of FPSs, and if $k\in\mathbb{N}$, then $\prod_{i\in
I}\mathbf{a}_{i}^{k}=\left(  \prod_{i\in I}\mathbf{a}_{i}\right)  ^{k}$. If
the $\mathbf{a}_{i}$ are moreover invertible, then this equality also holds
for all $k\in\mathbb{Z}$.}) allow us to work with infinite products almost as
comfortably as with finite ones. Of course, we need to check that our families
are multipliable, and occasionally verify that a few other technical
requirements are met (such as multipliability of subfamilies), but usually
such verifications are straightforward and easy and can be done in one's head.

Does this justify our manipulations in Subsection \ref{subsec.gf.prod.exa}? To
some extent. We need to be careful with the telescope principle, whose
infinite analogue is rather subtle and needs some qualifications. Here is an
example of how not to use the telescope principle:
\[
\dfrac{1}{2}\cdot\dfrac{2}{2}\cdot\dfrac{2}{2}\cdot\dfrac{2}{2}\cdot\cdots
\neq1.
\]
It is tempting to argue that the infinitely many $2$'s in these fractions
cancel each other out, and yet the $1$ that remains is not the right result.
See Exercise \ref{exe.fps.prods.telescope1} \textbf{(b)} for how an infinite
telescope principle should actually look like. Anyway, our computations in
Subsection \ref{subsec.gf.prod.exa} did not truly need the telescope
principle; they can just as well be made using more fundamental
rules\footnote{To wit, we can argue as follows: We have%
\[
\prod_{i\in\mathbb{N}}\left(  1+x^{2^{i}}\right)  =\prod_{i\in\mathbb{N}%
}\dfrac{1-x^{2^{i+1}}}{1-x^{2^{i}}}=\dfrac{\prod_{i\in\mathbb{N}}\left(
1-x^{2^{i+1}}\right)  }{\prod_{i\in\mathbb{N}}\left(  1-x^{2^{i}}\right)  },
\]
where the last step used Proposition \ref{prop.fps.div-mulable} and relied on
the fact that both families $\left(  1-x^{2^{i+1}}\right)  _{i\in\mathbb{N}}$
and $\left(  1-x^{2^{i}}\right)  _{i\in\mathbb{N}}$ are multipliable (this is
important, but very easy to check in this case) and that each FPS $1-x^{2^{i}%
}$ is invertible (because its constant term is $1$). However, splitting off
the factor for $i=0$ from the product $\prod_{i\in\mathbb{N}}\left(
1-x^{2^{i}}\right)  $, we obtain
\[
\prod_{i\in\mathbb{N}}\left(  1-x^{2^{i}}\right)  =\underbrace{\left(
1-x^{2^{0}}\right)  }_{=1-x^{1}=1-x}\cdot\underbrace{\prod_{i>0}\left(
1-x^{2^{i}}\right)  }_{\substack{=\prod_{i\in\mathbb{N}}\left(  1-x^{2^{i+1}%
}\right)  \\\text{(here, we substituted }i+1\\\text{for }i\text{ in the
product)}}}=\left(  1-x\right)  \cdot\prod_{i\in\mathbb{N}}\left(
1-x^{2^{i+1}}\right)  ,
\]
so that%
\[
\dfrac{\prod_{i\in\mathbb{N}}\left(  1-x^{2^{i+1}}\right)  }{\prod
_{i\in\mathbb{N}}\left(  1-x^{2^{i}}\right)  }=\dfrac{1}{1-x}.
\]
Hence, $\prod_{i\in\mathbb{N}}\left(  1+x^{2^{i}}\right)  =\dfrac{\prod
_{i\in\mathbb{N}}\left(  1-x^{2^{i+1}}\right)  }{\prod_{i\in\mathbb{N}}\left(
1-x^{2^{i}}\right)  }=\dfrac{1}{1-x}$. So we don't need the telescope
principle to justify this equality.}.

There is one more rule that we have used in Subsection
\ref{subsec.gf.prod.exa} and have not justified yet: the equality
(\ref{eq.fps.prod.binary.prod-inf}). We will justify it next.

